% =======================================================================
% Chapter V: The Thermal Instability of a Layer of Fluid Heated from Below
% Part 4: The Effect of Rotation and Magnetic Field
% Sections 49--54 + Bibliographical Notes
% Book pages 196--219
%
% Assembled from fragment files:
%   chapter_5_sec49_50.tex  (sections 49-50 now filled in)
%   chapter_5_sec51_52.tex  (complete)
%   chapter_5_sec53_bib.tex (complete)
% =======================================================================

\chapter{The Thermal Instability of a Layer of Fluid Heated from Below}
\label{ch:5}

\section*{4. THE EFFECT OF ROTATION AND MAGNETIC FIELD}

%% ==================================================================
%% Section 49
%% ==================================================================
\section{The like and the contrary effects of rotation and magnetic
  field on fluid behaviour}\label{sec:49}

In the last two chapters we have studied the effect of rotation and
magnetic field, acting separately, on the onset of thermal instability in
layers of fluid heated from below.  In some respects the effects are
remarkably alike: they both inhibit the onset of instability; and they
both elongate the cells which appear at marginal stability.  These effects
have likewise a common origin: in the Taylor--Proudman theorem in
the case of rotation and in its exact analogue in the case of a magnetic
field.  On these accounts one must not suppose that acting together,
rotation and magnetic field will reinforce each other; on the contrary,
in other respects they will tend to oppose each other.  Thus, we have
seen that viscosity facilitates the onset of instability when rotation is
present; and we have also seen that a magnetic field imparts to the fluid
certain aspects of viscosity.  Consequently, even though the two acting
separately inhibit the onset of instability, they will have conflicting
tendencies when acting together.  Again, there is a fundamental difference in the character of the motions which prevail when rotation alone
is present and when a magnetic field alone is present.  Rotation induces a
component of vorticity in the direction of~$\boldsymbol{\Omega}$, and the effects arising from
it are predominant; for large Taylor numbers it results in the streamlines
becoming closely wound spirals with motions principally confined to
planes transverse to~$\boldsymbol{\Omega}$.  But a magnetic field does not induce any similar
component of vorticity and there are no comparable effects; instead, for
large $Q$-numbers, the motions transverse to~$\mathbf{H}$ are much reduced and the
motions along the magnetic lines of force become predominant.

There is a further fact to remember: in liquid metals such as mercury,
instability sets in mostly as overstability when rotation is present, but
it sets in as stationary convection when a magnetic field is present.  For
all these reasons, the study of thermal instability in the presence of both
rotation and magnetic field is an instructive one.

%% ==================================================================
%% Section 50
%% ==================================================================
\section{The propagation of hydromagnetic waves in a rotating fluid}\label{sec:50}

In Chapters~III and~IV we saw that the propagation of waves in a
rotating fluid and in the presence of a magnetic field bear on the respective
problems in thermal instability.  We shall see that the propagation
of hydromagnetic waves in a rotating fluid has a similar bearing on the
problem to be investigated in this chapter.

Consider an incompressible fluid in a state of uniform rotation with
an angular velocity~$\boldsymbol{\Omega}$.  In a frame of reference rotating with the angular
velocity~$\boldsymbol{\Omega}$, the basic equations are (cf.\ equations III~(44) and IV~(10)):
\begin{align}
  \frac{\partial u_i}{\partial t}
  + u_j\,\frac{\partial u_i}{\partial x_j}
  - \frac{\mu H_j}{4\pi\rho}\,\frac{\partial H_i}{\partial x_j}
  \notag\\
  &= 2\epsilon_{ijk}\,u_j\,\Omega_k
    - \frac{\partial}{\partial x_i}\!\left(
        p + \frac{|\mathbf{H}|^{2}}{8\pi}
        - \tfrac{1}{2}|\boldsymbol{\Omega}\times\mathbf{r}|^{2} + V
    \right),
  \label{eq:5-1}
\end{align}
\begin{equation}
  \frac{\partial H_i}{\partial t}
  + u_j\,\frac{\partial H_i}{\partial x_j}
  - H_j\,\frac{\partial u_i}{\partial x_j}
  = \eta\,\nabla^{2}H_i,
  \label{eq:5-2}
\end{equation}
\begin{equation}
  \frac{\partial u_i}{\partial x_i} = 0,
  \quad\text{and}\quad
  \frac{\partial H_i}{\partial x_i} = 0.
  \label{eq:5-3}
\end{equation}

For an inviscid fluid of zero resistivity, the equations governing small
departures from an initial state in which a uniform magnetic field~$H_j$
prevails are
\begin{equation}
  \frac{\partial u_i}{\partial t}
  - \frac{\mu H_j}{4\pi\rho}\,\frac{\partial h_i}{\partial x_j}
  = -\frac{\partial}{\partial x_i}\,\delta\varpi
    + 2\epsilon_{ijk}\,u_j\,\Omega_k,
  \label{eq:5-4}
\end{equation}
\begin{equation}
  \frac{\partial h_i}{\partial t}
  = H_j\,\frac{\partial u_i}{\partial x_j},
  \quad
  \frac{\partial u_i}{\partial x_i} = 0,
  \quad\text{and}\quad
  \frac{\partial h_i}{\partial x_i} = 0,
  \label{eq:5-5}
\end{equation}
where~$h_i$ denotes the departure of the magnetic field from~$H_i$.

We now seek solutions of equations~\eqref{eq:5-4} and~\eqref{eq:5-5}
whose space-time dependence is given by
\begin{equation}
  e^{i(\sigma t + \mathbf{k}\cdot\mathbf{x})}.
  \label{eq:5-6}
\end{equation}

Equations~\eqref{eq:5-4} and~\eqref{eq:5-5} then give
\begin{equation}
  \sigma u_i
  - \frac{\mu(H_j k_j)}{4\pi\rho}\,h_i
  = -k_i\,\delta\varpi - 2i\epsilon_{ijk}\,u_j\,\Omega_k,
  \label{eq:5-7}
\end{equation}
\begin{equation}
  \sigma h_i = (H_j k_j)\,u_i,
  \label{eq:5-8}
\end{equation}
\begin{equation}
  u_i\,k_i = 0,
  \quad\text{and}\quad
  h_i\,k_i = 0.
  \label{eq:5-9}
\end{equation}

Equation~\eqref{eq:5-9} shows that \emph{the waves are transverse}.

Eliminating~$h_i$ between equations~\eqref{eq:5-7} and~\eqref{eq:5-8},
we obtain
\begin{equation}
  n u_i + k_i\,\delta\varpi + 2i\epsilon_{ijk}\,u_j\,\Omega_k = 0,
  \label{eq:5-10}
\end{equation}
where
\begin{equation}
  n = \frac{1}{\omega}(\omega^{2}-\omega_{A}^{2})
  \label{eq:5-11}
\end{equation}
and
\begin{equation}
  \omega_{A} = \left(\frac{\mu}{4\pi\rho}\right)^{\!\frac{1}{2}}
  (H_{j}\,k_{j}).
  \label{eq:5-12}
\end{equation}

It will be observed that equation~(10) is identical with equation~III~(61);
and the solution obtained there applies here.  Thus, by choosing the
orientation of the coordinate system such that the $z$-axis is the direction
of~$\mathbf{k}$ and the $(x,z)$-plane contains~$\boldsymbol{\Omega}$, we have
(cf.\ equations~III~(71) and~(73)),
\begin{equation}
  n = \pm\,2\Omega\cos\vartheta,
  \label{eq:5-13}
\end{equation}
\begin{equation}
  u_{x} = \pm\,i\,u_{y},\quad
  u_{z} = 0,\quad\text{and}\quad
  \delta\varpi = \mp\,\frac{2\Omega}{k}\,u_{z}\sin\vartheta,
  \label{eq:5-14}
\end{equation}
where $\vartheta$ is the inclination of the direction of wave propagation to
the direction of~$\boldsymbol{\Omega}$.  The corresponding amplitudes of the
perturbations in the magnetic field follow from equation~(8); we have
\begin{equation}
  h_{x} = \frac{k_{j}H_{j}}{\omega}\,u_{x},\quad
  h_{z} = \pm\,i\,h_{y},\quad\text{and}\quad
  h_{s} = 0.
  \label{eq:5-15}
\end{equation}

Now combining equations~(11) and~(13), we have
\begin{equation}
  \omega^{2} \mp (2\Omega\cos\vartheta)\,\omega - \omega_{A}^{2} = 0.
  \label{eq:5-16}
\end{equation}
Hence,
\begin{equation}
  \omega = \pm\bigl[\Omega\cos\vartheta
    \pm\sqrt{\Omega^{2}\cos^{2}\!\vartheta+\omega_{A}^{2}}\,\bigr].
  \label{eq:5-17}
\end{equation}
If $\omega_{1}$ and $\omega_{2}$ are the two roots of equation~\eqref{eq:5-16},
then
\begin{equation}
  \omega_{1}+\omega_{2} = \pm\,2\Omega\cos\vartheta,\qquad
  \omega_{1}\,\omega_{2} = -\omega_{A}^{2}.
  \label{eq:5-18}
\end{equation}

The velocity of propagation of the waves is given by
\begin{equation}
  V = \frac{\omega}{k}
    = \pm\,\frac{1}{k}
    \bigl[\Omega\cos\vartheta
      \pm\sqrt{\Omega^{2}\cos^{2}\!\vartheta+\omega_{A}^{2}}\,\bigr].
  \label{eq:5-19}
\end{equation}
This is the phase velocity.  The corresponding group velocity is given
by (cf.\ equations~III~(75) and~IV~(64))
\begin{equation}
  \frac{\partial\omega}{\partial\mathbf{k}}
  = \pm\left\{
    \frac{\mathbf{k}\times(\boldsymbol{\Omega}\times\mathbf{k})}{k^{3}}
    \pm\frac{\Omega\cos\vartheta\,
      [\mathbf{k}\times(\boldsymbol{\Omega}\times\mathbf{k})]\,k^{-3}
      +\omega_{A}\,\mathbf{H}(\mu/4\pi\rho)^{\frac{1}{2}}}%
    {\sqrt{\Omega^{2}\cos^{2}\!\vartheta+\omega_{A}^{2}}}
  \right\}.
  \label{eq:5-20}
\end{equation}

%% ==================================================================
%% Section 51
%% ==================================================================
\section{The perturbation equations}\label{sec:51}

% [Continues from previous page]
\noindent
\ldots\ last chapters, that the basic perturbation equations are (cf.\ equations
III~(78)--(80) and IV~(87)--(92)):
\begin{equation}
  \frac{\partial u_i}{\partial t}
  = -\frac{\partial}{\partial x_i}(\delta\varpi)
  + g\alpha\Theta\,\delta_{i3}
  + \nu\nabla^{2}u_i
  + 2\epsilon_{ijk}\,u_j\,\Omega_k
  + \frac{\mu H_j}{4\pi\rho}\,\frac{\partial h_i}{\partial x_j},
  \label{eq:5-21}
\end{equation}
\begin{equation}
  \frac{\partial h_i}{\partial t}
  = H_j\,\frac{\partial u_i}{\partial x_j}+\eta\nabla^{2}h_i,
  \label{eq:5-22}
\end{equation}
\begin{equation}
  \frac{\partial\Theta}{\partial t}
  = \beta w + \kappa\nabla^{2}\Theta,
  \label{eq:5-23}
\end{equation}
\begin{equation}
  \frac{\partial u_i}{\partial x_i} = 0, \quad\text{and}\quad
  \frac{\partial h_i}{\partial x_i} = 0.
  \label{eq:5-24}
\end{equation}

By taking the curl of equation~\eqref{eq:5-21}, we can eliminate the term
in~$\delta\varpi$; we thus obtain (cf.\ equations III~(82) and IV~(93))
\begin{equation}
  \frac{\partial\omega_i}{\partial t}
  = g\alpha\,\epsilon_{ijk}\,\frac{\partial\Theta}{\partial x_j}\,\lambda_k
  + \nu\nabla^{2}\omega_i
  + 2\Omega_j\,\frac{\partial u_i}{\partial x_j}
  + \frac{\mu H_j}{4\pi\rho}\,\frac{\partial v_i}{\partial x_j},
  \label{eq:5-25}
\end{equation}
where
\begin{equation}
  \boldsymbol{\omega} = \operatorname{curl}\mathbf{u}
  \quad\text{and}\quad
  \mathbf{v} = \operatorname{curl}\mathbf{h}.
  \label{eq:5-26}
\end{equation}

Taking the curl of equation~\eqref{eq:5-25} once again, we obtain
(cf.\ equations III~(83) and IV~(95))
\begin{equation}
  \frac{\partial}{\partial t}\nabla^{2}u_i
  = g\alpha\!\left(
      \lambda_i\,\nabla^{2}\Theta
      - \lambda_j\,\frac{\partial^{2}\Theta}{\partial x_i\,\partial x_j}
    \right)
  + \nu\nabla^{4}u_i
  - 2\Omega_j\,\frac{\partial\omega_i}{\partial x_j}
  + \frac{\mu H_j}{4\pi\rho}\,\frac{\partial}{\partial x_j}\,\nabla^{2}h_i.
  \label{eq:5-27}
\end{equation}

Now multiplying equations~\eqref{eq:5-25} and~\eqref{eq:5-27} by~$\delta_{i3}$
we get
\begin{equation}
  \frac{\partial\zeta}{\partial t}
  = \nu\nabla^{2}\zeta
  + 2\Omega\,\frac{\partial w}{\partial x_3}
  + \frac{\mu_H}{4\pi}\,\frac{\partial\xi}{\partial x_3},
  \label{eq:5-28}
\end{equation}
and
\begin{equation}
  \frac{\partial}{\partial t}\nabla^{2}w
  = g\alpha\!\left(
      \frac{\partial^{2}\Theta}{\partial x_1^{2}}
      + \frac{\partial^{2}\Theta}{\partial x_2^{2}}
    \right)
  + \nu\nabla^{4}w
  - 2\Omega\,\frac{\partial\zeta}{\partial x_3}
  + \frac{\mu_H}{4\pi}\,\frac{\partial}{\partial x_3}\,\nabla^{2}h_3,
  \label{eq:5-29}
\end{equation}
where $\zeta$ and~$\xi/4\pi$ are the $z$-components of the vorticity and the
current density, respectively.  Equations~\eqref{eq:5-28} and~\eqref{eq:5-29}
replace equations~IV~(99) and~(100); the remaining equations IV~(96)--(98) are
unaffected.

We shall restrict our discussion of this problem to the case when
$\mathbf{g}$, $\mathbf{H}$, and~$\boldsymbol{\Omega}$ act in the same
direction.  In this case the relevant equations are:
\begin{equation}
  \frac{\partial\Theta}{\partial t}
  = \kappa\nabla^{2}\Theta + \beta w,
  \label{eq:5-30}
\end{equation}
\begin{equation}
  \frac{\partial h_3}{\partial t}
  = \eta\nabla^{2}h_3 + H\,\frac{\partial w}{\partial z},
  \label{eq:5-31}
\end{equation}
\begin{equation}
  \frac{\partial\xi}{\partial t}
  = \eta\nabla^{2}\xi + H\,\frac{\partial\zeta}{\partial z},
  \label{eq:5-32}
\end{equation}
\begin{equation}
  \frac{\partial\zeta}{\partial t}
  = \nu\nabla^{2}\zeta
  + 2\Omega\,\frac{\partial w}{\partial z}
  + \frac{\mu H}{4\pi\rho}\,\frac{\partial\xi}{\partial z},
  \label{eq:5-33}
\end{equation}
and
\begin{equation}
  \frac{\partial}{\partial t}(\nabla^{2}w)
  = g\alpha\!\left(
      \frac{\partial^{2}\Theta}{\partial x^{2}}
      + \frac{\partial^{2}\Theta}{\partial y^{2}}
    \right)
  + \nu\nabla^{4}w
  - 2\Omega\,\frac{\partial\zeta}{\partial z}
  + \frac{\mu H}{4\pi\rho}\,\frac{\partial}{\partial z}\,\nabla^{2}h_3.
  \label{eq:5-34}
\end{equation}

Analysing the disturbances into two-dimensional waves, and considering
disturbances characterized by a particular wave number, we find that
equations~\eqref{eq:5-30}--\eqref{eq:5-34} give
\begin{equation}
  (D^{2}-a^{2}-p_1\,\sigma)\,\Theta
  = -\left(\frac{\beta d^{2}}{\kappa}\right)W,
  \label{eq:5-35}
\end{equation}
\begin{equation}
  (D^{2}-a^{2}-p_2\,\sigma)\,K
  = -\left(\frac{Hd}{\eta}\right)DW,
  \label{eq:5-36}
\end{equation}
\begin{equation}
  (D^{2}-a^{2}-p_2\,\sigma)\,X
  = -\left(\frac{Hd}{\eta}\right)DZ,
  \label{eq:5-37}
\end{equation}
\begin{equation}
  (D^{2}-a^{2}-\sigma)\,Z
  = -\left(\frac{2\Omega d^{2}}{\nu}\right)DW
  - \left(\frac{\mu Hd}{4\pi\rho\nu}\right)DX,
  \label{eq:5-38}
\end{equation}
and
\begin{multline}
  (D^{2}-a^{2})
  (D^{2}-a^{2}-\sigma)\,W
  + \left(\frac{\mu Hd}{4\pi\rho\nu}\right)
  D(D^{2}-a^{2})\,K \\
  - \left(\frac{2\Omega d^{2}}{\nu}\right)DZ
  \equiv \left(\frac{g\alpha d^{2}}{\nu}\right)a^{2}\Theta,
  \label{eq:5-39}
\end{multline}
where the notation is the same as in equations~IV~(118)--(123).

Eliminating $X$ between equations~\eqref{eq:5-37} and~\eqref{eq:5-38},
and similarly~$K$ between equations~\eqref{eq:5-36}
and~\eqref{eq:5-39}, we obtain
\begin{multline}
  \bigl[(D^{2}-a^{2}-\sigma)(D^{2}-a^{2}-p_2\,\sigma)
    -QD^{2}\bigr]Z \\
  = -\left(\frac{2\Omega d^{2}}{\nu}\right)
  D(D^{2}-a^{2}-p_2\,\sigma)\,W
  \label{eq:5-40}
\end{multline}
and
\begin{multline}
  \bigl[D^{2}-a^{2}\bigr]
  \bigl[(D^{2}-a^{2}-\sigma)(D^{2}-a^{2}-p_2\,\sigma)
    -QD^{2}\bigr]W \\
  -\left(\frac{2\Omega d^{2}}{\nu}\right)
  D(D^{2}-a^{2}-p_2\,\sigma)\,Z
  = \left(\frac{g\alpha d^{2}}{\nu}\right)
  a^{2}(D^{2}-a^{2}-p_2\,\sigma)\,\Theta.
  \label{eq:5-41}
\end{multline}

Eliminating $Z$ between these last two equations, we have
\begin{multline}
  \bigl\{(D^{2}\!-\!a^{2})
    \bigl[(D^{2}\!-\!a^{2}\!-\!\sigma)(D^{2}\!-\!a^{2}\!-\!p_2\,\sigma)
      -QD^{2}\bigr]^{2} \\
  +\,T D^{2}(D^{2}\!-\!a^{2}\!-\!p_2\,\sigma)^{2}\bigr\}W \\
  = \left(\frac{g\alpha d^{2}}{\nu}\right)
  a^{2}\bigl[(D^{2}\!-\!a^{2}\!-\!\sigma)(D^{2}\!-\!a^{2}\!-\!p_2\,\sigma)
    -QD^{2}\bigr]
  (D^{2}\!-\!a^{2}\!-\!p_2\,\sigma)\,\Theta.
  \label{eq:5-42}
\end{multline}

Finally, eliminating~$\Theta$ between equations~\eqref{eq:5-35}
and~\eqref{eq:5-42} we obtain
\begin{multline}
  (D^{2}\!-\!a^{2}\!-\!p_1\,\sigma)
  \bigl\{(D^{2}\!-\!a^{2})
    \bigl[(D^{2}\!-\!a^{2}\!-\!\sigma)(D^{2}\!-\!a^{2}\!-\!p_2\,\sigma)
      -QD^{2}\bigr]^{2} \\
  +\,T D^{2}(D^{2}\!-\!a^{2}\!-\!p_2\,\sigma)^{2}\bigr\}W \\
  = -\mathit{Ra}^{2}
    \bigl[(D^{2}\!-\!a^{2}\!-\!\sigma)(D^{2}\!-\!a^{2}\!-\!p_2\,\sigma)
    -QD^{2}\bigr]
    (D^{2}\!-\!a^{2}\!-\!p_2\,\sigma)\,W.
  \label{eq:5-43}
\end{multline}

Solutions of the foregoing equations must be sought which satisfy the
boundary conditions given in equations~II~(96) and~(96) and IV~(124)
and~(125).

It may be recalled here that all equations following~\eqref{eq:5-21} apply only if
$\mathbf{g}$, $\mathbf{H}$, and~$\boldsymbol{\Omega}$ are parallel.  The
generalization to the case when $\mathbf{g}$, $\mathbf{H}$,
and~$\boldsymbol{\Omega}$ are coplanar is immediate.  By restricting oneself
to the onset of convection in the form of infinitely extended rolls (along
directions parallel to the plane containing $\mathbf{g}$, $\mathbf{H}$,
and~$\boldsymbol{\Omega}$), instead of cells, one can readily verify that the
same set of equations applies to this more general case if we interpret $H$
and~$\Omega$ in all the equations to mean the components of~$\mathbf{H}$
and~$\boldsymbol{\Omega}$ in the direction of the vertical.  The restriction
to the onset of convection in the form of rolls when $\mathbf{g}$,
$\mathbf{H}$, and~$\boldsymbol{\Omega}$ are coplanar can very likely be
justified in the way in which the same restriction was justified in
Chapter~IV, \S\,47 for the case when $\mathbf{H}$ alone was acting.
However, when $\mathbf{g}$, $\mathbf{H}$, and~$\boldsymbol{\Omega}$ are not
coplanar, the solution of the problem will require an explicit minimizing of
the Rayleigh number as a function of the two wave numbers characterizing the
disturbance in two directions at right angles to one another and in the
horizontal plane.  But it is unlikely that the solution will exhibit any
essentially new feature not already disclosed by the solution for the case
when $\mathbf{g}$, $\mathbf{H}$, and~$\boldsymbol{\Omega}$ are parallel.

%% ==================================================================
%% Section 52
%% ==================================================================
\section{The case when instability sets in as stationary convection}\label{sec:52}

We shall first consider the case when instability sets in as ordinary
convection.  The case when it sets in as overstability is considered in the
following section.

When instability sets in as ordinary convection, the marginal state will be
characterized by $\sigma = 0$; and the basic equations are:
\begin{equation}
  (D^{2}-a^{2})\,\Theta
  = -\left(\frac{\beta d^{2}}{\kappa}\right)W,
  \label{eq:5-44}
\end{equation}
\begin{equation}
  \bigl[(D^{2}-a^{2})^{2}-QD^{2}\bigr]Z
  = -\left(\frac{2\Omega d^{2}}{\nu}\right)D(D^{2}-a^{2})\,W,
  \label{eq:5-45}
\end{equation}
\begin{equation}
  (D^{2}-a^{2})\,K
  = -\left(\frac{Hd}{\eta}\right)DW,
  \label{eq:5-46}
\end{equation}
\begin{equation}
  (D^{2}-a^{2})\,X
  = -\left(\frac{Hd}{\eta}\right)DZ,
  \label{eq:5-47}
\end{equation}
and
\begin{multline}
  (D^{2}-a^{2})
  \bigl\{\bigl[(D^{2}-a^{2})^{2}-QD^{2}\bigr]^{2}
  +T D^{2}(D^{2}-a^{2})^{2}\bigr\}W \\
  = -\mathit{Ra}^{2}\bigl\{(D^{2}-a^{2})^{2}-QD^{2}\bigr\}W.
  \label{eq:5-48}
\end{multline}

According to equations~II~(95) and~(96) and IV~(124) and~(125), the
boundary conditions with respect to which equations~\eqref{eq:5-44}--\eqref{eq:5-48}
must be solved are:
\begin{equation}
  W = 0 \quad\text{and}\quad \Theta = 0 \quad\text{for } z = 0 \text{ and } 1,
  \label{eq:5-49}
\end{equation}
\begin{equation}
  \left.
  \begin{aligned}
    &\textit{either}\quad DW = 0 \quad\text{and}\quad Z = 0
      \quad\text{(on a rigid boundary)} \\
    &\textit{or}\quad D^{2}W = 0 \quad\text{and}\quad DZ = 0
      \quad\text{(on a free boundary)}
  \end{aligned}
  \right\}
  \label{eq:5-50}
\end{equation}
and
\begin{equation}
  \left.
  \begin{aligned}
    &\textit{either}\quad DX = 0 \;\text{and}\; K = 0
      \;\text{(on a perfectly conducting boundary)}, \\
    &\textit{or}\quad X = 0
      \;\text{(on a boundary adjoining a non-conducting medium)}.
  \end{aligned}
  \right\}
  \label{eq:5-51}
\end{equation}

The foregoing equations and boundary conditions constitute a characteristic
value problem in a differential equation of order ten.  In spite of a
variational formulation which is possible (see Appendix~II), the problem is
clearly of considerable complexity.  For this reason, we shall restrict
ourselves to the case when both the bounding surfaces are free and the medium
adjoining the fluid is non-conducting.  While this is admittedly an artificial
case to consider, we do not expect to lose any of the essential features of the
problem on this account; for, we have seen in Chapters~III and~IV that the
solutions obtained for different sets of boundary conditions all exhibit the
same dependence on the parameters of the problem such as $\mathrm{p}$, $T$,
and~$Q$.

\subsection*{(a) \textit{The solution for the case of two free boundaries}}

As in Chapter~III, \S\,27\,(a) and Chapter~IV, \S\,44\,(a), it can be shown
that for this case, the proper solutions of
equations~\eqref{eq:5-44}--\eqref{eq:5-48} are given by
\begin{equation}
  W = W_0\sin n\pi z,
  \label{eq:5-52}
\end{equation}
\begin{equation}
  \Theta = \frac{\beta}{\kappa}\,
  \frac{d^{2}}{n^{2}\pi^{2}+a^{2}}\,W_0\sin n\pi z,
  \label{eq:5-53}
\end{equation}
\begin{equation}
  Z = \frac{2\Omega d}{\nu}\,
  \frac{n\pi(n^{2}\pi^{2}+a^{2})}{(n^{2}\pi^{2}+a^{2})^{2}+Qn^{2}\pi^{2}}\,
  W_0\cos n\pi z,
  \label{eq:5-54}
\end{equation}
\begin{equation}
  K = \frac{Hd}{\eta}\,
  \frac{n\pi}{n^{2}\pi^{2}+a^{2}}\,
  W_0\cos n\pi z + \text{constant} \times \cosh az,
  \label{eq:5-55}
\end{equation}
and
\begin{equation}
  X = \frac{2\Omega d}{\nu}\,\frac{Hd}{\eta}\,
  \frac{n^{2}\pi^{2}}{(n^{2}\pi^{2}+a^{2})^{2}+Q n^{2}\pi^{2}}\,
  W_0\sin n\pi z,
  \label{eq:5-56}
\end{equation}
where $n$ is an integer.  The corresponding characteristic equation is
obtained by substituting the solution~\eqref{eq:5-52} for~$W$ in
equation~\eqref{eq:5-48}.  We find
\begin{equation}
  R = \frac{%
    (n^{2}\pi^{2}+a^{2})
    \bigl\{\bigl[(n^{2}\pi^{2}+a^{2})^{2}+Qn^{2}\pi^{2}\bigr]^{2}
    +Tn^{2}\pi^{2}(n^{2}\pi^{2}+a^{2})\bigr\}}{%
    a^{2}\bigl[(n^{2}\pi^{2}+a^{2})^{2}+Qn^{2}\pi^{2}\bigr]}.
  \label{eq:5-57}
\end{equation}

Letting $a^{2} = \pi^{2}x$, we can rewrite equation~\eqref{eq:5-57} in
the form
\begin{equation}
  R = \pi^{4}\,\frac{n^{2}+x}{x}
  \left[
    (n^{2}+x)^{2}
    + \frac{Qn^{2}}{\pi^{2}}
  \right]
  + \frac{T}{x}\,
  \frac{1}{1/n^{2}+Q/\pi^{2}(n^{2}+x)^{2}}.
  \label{eq:5-58}
\end{equation}

From this equation it follows that instability first sets in for the lowest
mode $n = 1$.  The corresponding expression for~$R$ is
\begin{equation}
  R = \pi^{4}\,\frac{%
    (1+x)
    \bigl\{\bigl[(1+x)^{2}+Q_1\bigr]^{2}
    +T_1(1+x)\bigr\}}{%
    x\bigl[(1+x)^{2}+Q_1\bigr]},
  \label{eq:5-59}
\end{equation}
where
\begin{equation}
  x = \frac{a^{2}}{\pi^{2}}, \quad
  Q_1 = \frac{Q}{\pi^{2}}, \quad\text{and}\quad
  T_1 = \frac{T}{\pi^{4}}.
  \label{eq:5-60}
\end{equation}

As a function of~$x$, $R$ given by equation~\eqref{eq:5-59} attains its
extremal value when
\begin{equation}
  2x^{3}+3x^{2}-1
  = Q_1 + T_1\,
  \frac{(1+x)^{4}-Q_1(x^{2}-1)}{%
    \bigl[(1+x)^{2}+Q_1\bigr]^{2}}.
  \label{eq:5-61}
\end{equation}

\figurePlaceholder{47}{The critical Rayleigh number $R_c$ for the cases of
ordinary cellular convection (solid line) and overstability (for
$\mathrm{p} = 0.025$) (broken line) as a function of
$Q_1\,(= Q/\pi^{2})$ for various assigned values of
$T_1\,(= T/\pi^{4})$.  The curves are labeled by the values of~$T_1$ to which
they refer.  For a given value of~$T_1$, instability will set in as
overstability for all values of~$Q_1$ less than that at the point of
intersection of the corresponding full-line and dashed curves; for all larger
values of~$Q_1$, it will set in as ordinary convection.}

However, this equation is not very useful for determining the critical Rayleigh
numbers for assigned values of~$Q_1$ and~$T_1$.  It is more convenient to
evaluate~$R$ directly as a function of~$x$ (in accordance with
equation~\eqref{eq:5-59}) and locate the minimum numerically.  The critical
numbers listed in Table~XVIII and illustrated in Figs.~47 and~48 were
determined in this fashion.

It is apparent from Table~XVIII and Figs.~47 and~48 that the solution

% ----------------------------------------------------------------
% TABLE XVIII
% ----------------------------------------------------------------
\begin{table}[p]
\centering
\caption{The critical Rayleigh numbers and the wave numbers of the associated
disturbance for the onset of instability as stationary convection for various
values of $Q_1\,(= Q/\pi^{2})$ and $T_1\,(= T/\pi^{4})$.}
\label{tab:XVIII}

\medskip
\small
\begin{tabular}{r@{\quad}r@{\quad}r|r@{\quad}r@{\quad}r}
\hline
\multicolumn{3}{c|}{$T_1 = 0$} & \multicolumn{3}{c}{$T_1 = 10$} \\
$Q_1$ & $a$ & $R_c$ & $Q_1$ & $a$ & $R_c$ \\
\hline
 0     & $3{\cdot}70$  &           & 0     & $4{\cdot}31$  &           \\
 10    & $3{\cdot}97$  & $6{\cdot}95\times10^{3}$ & 10    & $4{\cdot}56$  &           \\
 40    &               &           & 40    &               &           \\
 100   & $5{\cdot}67$  &           & 100   & $5{\cdot}96$  &           \\
 1{,}000 & $8{\cdot}59$ &          & 1{,}000 & $8{\cdot}65$ &          \\
 10{,}000 & $12{\cdot}6$ &$1{\cdot}06\times10^{6}$& 10{,}000 & $12{\cdot}6$  & \\
 50{,}000 & $16{\cdot}8$ &$5{\cdot}09\times10^{6}$& 50{,}000 & $16{\cdot}8$  & \\
 100{,}000 & $18{\cdot}94$&$1{\cdot}01\times10^{7}$& 100{,}000 & $18{\cdot}94$ & \\
\hline
\end{tabular}

\bigskip
\begin{tabular}{r@{\quad}r@{\quad}r|r@{\quad}r@{\quad}r}
\hline
\multicolumn{3}{c|}{$T_1 = 50$} & \multicolumn{3}{c}{$T_1 = 100$} \\
$Q_1$ & $a$ & $R_c$ & $Q_1$ & $a$ & $R_c$ \\
\hline
 0    & $5{\cdot}04$  & $3{\cdot}88\times10^{4}$ & 0    & $5{\cdot}73$ &$1{\cdot}03\times10^{5}$\\
 10   & $5{\cdot}30$  &                          & 10   & $5{\cdot}92$ &           \\
 40   & $5{\cdot}96$  & $6{\cdot}43\times10^{4}$ & 40   & $6{\cdot}45$ &           \\
 100  & $6{\cdot}95$  &                          & 100  & $7{\cdot}20$ &           \\
 1{,}000 & $10{\cdot}13$ & $1{\cdot}14\times10^{5}$& 1{,}000 &             &           \\
 10{,}000 & $12{\cdot}6$ &                        & 10{,}000 & $12{\cdot}6$ &          \\
 50{,}000 & $16{\cdot}8$ &                        & 50{,}000 & $16{\cdot}8$ &          \\
 100{,}000 & $18{\cdot}94$ &                      & 100{,}000 & $18{\cdot}94$ &        \\
\hline
\end{tabular}

\bigskip
\begin{tabular}{r@{\quad}r@{\quad}r|r@{\quad}r@{\quad}r}
\hline
\multicolumn{3}{c|}{$T_1 = 500$} & \multicolumn{3}{c}{$T_1 = 1{,}000$} \\
$Q_1$ & $a$ & $R_c$ & $Q_1$ & $a$ & $R_c$ \\
\hline
 0    & $7{\cdot}90$  & $5{\cdot}03\times10^{5}$ & 0    & $9{\cdot}05$  & $2{\cdot}05\times10^{6}$\\
 10   & $7{\cdot}98$  &                          & 20   & $9{\cdot}14$  &           \\
 20   & $8{\cdot}08$  & $1{\cdot}73\times10^{5}$ & 40   & $9{\cdot}28$  &           \\
 40   & $8{\cdot}39$  &                          & 50   & $9{\cdot}36$  &           \\
 50   & $8{\cdot}60$  & $2{\cdot}05\times10^{5}$ & 80   & $9{\cdot}60$  &           \\
 80   & $9{\cdot}19$  &                          & 100  & $9{\cdot}80$  &           \\
 100  & $9{\cdot}52$  & $3{\cdot}45\times10^{5}$ & 500  & $11{\cdot}78$ &           \\
 500  & $11{\cdot}46$ & $6{\cdot}93\times10^{5}$ & 1{,}000 & $12{\cdot}94$ &        \\
 1{,}000 & $12{\cdot}77$ & $1{\cdot}08\times10^{6}$& 10{,}000 & $18{\cdot}0$ &      \\
 10{,}000 & $17{\cdot}8$ & $8{\cdot}65\times10^{6}$& 50{,}000 &             &       \\
 50{,}000 & $16{\cdot}8$ &                        & 100{,}000 & $18{\cdot}94$ &      \\
 100{,}000 & $18{\cdot}94$ & $1{\cdot}01\times10^{7}$ & & & \\
\hline
\end{tabular}

\bigskip
\begin{tabular}{r@{\quad}r@{\quad}r|r@{\quad}r@{\quad}r}
\hline
\multicolumn{3}{c|}{$T_1 = 1{,}500$} & \multicolumn{3}{c}{$T_1 = 2{,}000$} \\
$Q_1$ & $a$ & $R_c$ & $Q_1$ & $a$ & $R_c$ \\
\hline
 0    & $9{\cdot}87$  & $4{\cdot}05\times10^{6}$ & 0    & $10{\cdot}5$ & $6{\cdot}78\times10^{6}$\\
 10   & $9{\cdot}90$  &                          & 10   & $10{\cdot}5$ &           \\
 20   & $9{\cdot}95$  &                          & 20   & $10{\cdot}5$ &           \\
 30   & $10{\cdot}0$  & $4{\cdot}53\times10^{6}$ & 30   & $10{\cdot}5$ &           \\
 40   & $10{\cdot}1$  &                          & 40   & $10{\cdot}6$ &           \\
 50   & $10{\cdot}1$  & $5{\cdot}14\times10^{6}$ & 50   & $10{\cdot}6$ & $8{\cdot}17\times10^{6}$\\
 80   & $10{\cdot}4$  &                          & 80   & $10{\cdot}8$ &           \\
 100  & $10{\cdot}6$  & $6{\cdot}55\times10^{6}$ & 100  & $10{\cdot}9$ &           \\
 500  & $12{\cdot}20$ &                          & 500  & $12{\cdot}57$ &          \\
 1{,}000 & $13{\cdot}25$ &                       & 1{,}000 & $13{\cdot}60$ &       \\
 10{,}000 & $18{\cdot}2$ &$8{\cdot}65\times10^{6}$& 10{,}000 & $18{\cdot}3$ &     \\
 50{,}000 & $16{\cdot}94$ &                      & 50{,}000 &              &       \\
 100{,}000 & $18{\cdot}94$ & $1{\cdot}02\times10^{7}$ & 100{,}000 & $18{\cdot}94$ & \\
\hline
\end{tabular}

\bigskip
\begin{tabular}{r@{\quad}r@{\quad}r|r@{\quad}r@{\quad}r}
\hline
\multicolumn{3}{c|}{$T_1 = 5{,}000$} & \multicolumn{3}{c}{$T_1 = 10{,}000$} \\
$Q_1$ & $a$ & $R_c$ & $Q_1$ & $a$ & $R_c$ \\
\hline
 10   & $12{\cdot}02$ & $4{\cdot}00\times10^{7}$ & 10   & $12{\cdot}98$ & $6{\cdot}09\times10^{7}$\\
 20   & $9{\cdot}95$  &                          & 20   &              & $8{\cdot}88\times10^{7}$\\
 40   & $9{\cdot}82$  & $4{\cdot}00\times10^{7}$ & 40   & $11{\cdot}09$ &           \\
 50   & $7{\cdot}00$  & $3{\cdot}93\times10^{7}$ & 50   & $8{\cdot}69$  &           \\
 80   & $4{\cdot}94$  &                          & 80   & $8{\cdot}68$  &           \\
 100  & $4{\cdot}55$  & $3{\cdot}77\times10^{7}$ & 100  & $3{\cdot}77$  &           \\
 200  & $4{\cdot}13$  &                          & 200  & $4{\cdot}73$  & $7{\cdot}21\times10^{7}$\\
 500  & $5{\cdot}66$  &                          & 500  & $5{\cdot}31$  &           \\
 1{,}000 & $7{\cdot}22$ & $6{\cdot}89\times10^{7}$& 1{,}000 & $6{\cdot}51$ &       \\
 10{,}000 & $12{\cdot}6$ & $1{\cdot}21\times10^{8}$& 10{,}000 &            &       \\
 50{,}000 &             &                        & 50{,}000 &              &       \\
 100{,}000 & $18{\cdot}94$ &$1{\cdot}03\times10^{7}$& 100{,}000 & $18{\cdot}94$ & \\
\hline
\end{tabular}

\bigskip
\begin{tabular}{r@{\quad}r@{\quad}r|r@{\quad}r@{\quad}r}
\hline
\multicolumn{3}{c|}{$T_1 = 50{,}000$} & \multicolumn{3}{c}{$T_1 = 100{,}000$} \\
$Q_1$ & $a$ & $R_c$ & $Q_1$ & $a$ & $R_c$ \\
\hline
 10   & $19{\cdot}3$ & $1{\cdot}86\times10^{8}$ & 10   & $18{\cdot}9$ & $4{\cdot}15\times10^{8}$\\
 20   & $19{\cdot}1$ &                          & 20   & $18{\cdot}8$ &           \\
 40   & $19{\cdot}1$ & $1{\cdot}81\times10^{8}$ & 40   & $18{\cdot}4$ &           \\
 50   & $14{\cdot}9$  & $1{\cdot}86\times10^{8}$& 50   & $18{\cdot}2$ & $3{\cdot}96\times10^{8}$\\
 80   & $14{\cdot}8$  &                         & 80   &              &           \\
 100  & $14{\cdot}54$ & $3{\cdot}86\times10^{8}$& 100  & $3{\cdot}54$ & $1{\cdot}37\times10^{8}$\\
 200  & $3{\cdot}97$ & $5{\cdot}94\times10^{7}$ & 200  & $3{\cdot}48$ &           \\
 500  & $3{\cdot}71$  &                         & 500  & $4{\cdot}33$ &           \\
 1{,}000 & $3{\cdot}09$ &                       & 1{,}000 & $3{\cdot}97$ & $1{\cdot}37\times10^{8}$\\
 4{,}000 &             &                        & 4{,}000 &              &           \\
 10{,}000 & $10{\cdot}5$ &$1{\cdot}08\times10^{8}$& 10{,}000 &            &          \\
 30{,}000 &             &                        & 30{,}000 &              &         \\
 50{,}000 & $16{\cdot}6$ &                       & 50{,}000 &              &         \\
 100{,}000 & $18{\cdot}94$ & $1{\cdot}03\times10^{7}$ & 100{,}000 & $18{\cdot}94$ & \\
\hline
\end{tabular}

\bigskip
\begin{tabular}{r@{\quad}r@{\quad}r|r@{\quad}r@{\quad}r}
\hline
\multicolumn{3}{c|}{$T_1 = 10^{5}$} & \multicolumn{3}{c}{$T_1 = 10^{6}$} \\
$Q_1$ & $a$ & $R_c$ & $Q_1$ & $a$ & $R_c$ \\
\hline
 10   & $47{\cdot}8$ & $1{\cdot}85\times10^{8}$ & 10   & $41{\cdot}01$ &           \\
 50   & $47{\cdot}2$ &                          & 50   & $41{\cdot}01$ &           \\
 80   & $47{\cdot}2$ &                          & 80   & $35{\cdot}43$ &           \\
 100  & $47{\cdot}3$ & $1{\cdot}44\times10^{8}$ & 100  &              &           \\
 200  & $37{\cdot}1$ &                          & 200  &              &           \\
 300  & $3{\cdot}35$ &                          & 300  & $3{\cdot}19$ &           \\
 500  & $3{\cdot}12$ &                          & 500  & $3{\cdot}10$ &           \\
 1{,}000 & $3{\cdot}35$ &$9{\cdot}03\times10^{7}$& 1{,}000 &            &          \\
 5{,}000 & $5{\cdot}13$ &                       & 5{,}000 &              &         \\
 10{,}000 & $9{\cdot}33$ &$1{\cdot}06\times10^{8}$& 10{,}000 &           &          \\
 100{,}000 & $17{\cdot}9$ &                     & 100{,}000 &             &          \\
 1{,}000{,}000 & $18{\cdot}94$ &                & 1{,}000{,}000 & $18{\cdot}94$ &   \\
\hline
\end{tabular}

\bigskip
\begin{tabular}{r@{\quad}r@{\quad}r|r@{\quad}r@{\quad}r}
\hline
\multicolumn{3}{c|}{$T_1 = 10^{7}$} & \multicolumn{3}{c}{$T_1 = 10^{8}$} \\
$Q_1$ & $a$ & $R_c$ & $Q_1$ & $a$ & $R_c$ \\
\hline
 10   & $60{\cdot}96$ & $3{\cdot}97\times10^{9}$ & 50   & $88{\cdot}47$ &           \\
 50   & $60{\cdot}23$ &                          & 100  & $88{\cdot}45$ &           \\
 100  & $60{\cdot}14$ & $3{\cdot}97\times10^{9}$ & 200  & $88{\cdot}43$ &           \\
 500  & $60{\cdot}86$ & $1{\cdot}38\times10^{9}$ & 500  & $88{\cdot}34$ &$1{\cdot}93\times10^{9}$\\
 800  & $60{\cdot}86$ & $4{\cdot}86\times10^{8}$ & 1{,}000 & $88{\cdot}24$ &       \\
 1{,}000 & $3{\cdot}47$ &                       & 3{,}000 & $3{\cdot}96$ &          \\
 5{,}000 & $3{\cdot}14$ &$3{\cdot}96\times10^{8}$& 5{,}000 & $3{\cdot}15$ &        \\
 10{,}000 & $3{\cdot}74$ &                      & 10{,}000 & $3{\cdot}36$ &$1{\cdot}03\times10^{8}$\\
 100{,}000 & $17{\cdot}4$ & $1{\cdot}03\times10^{8}$ & 100{,}000 &         &        \\
 1{,}000{,}000 & $18{\cdot}94$ &                & 1{,}000{,}000 &$18{\cdot}94$&     \\
\hline
\end{tabular}

\bigskip
\begin{tabular}{r@{\quad}r@{\quad}r}
\hline
\multicolumn{3}{c}{$T_1 = 10^{9}$} \\
$Q_1$ & $a$ & $R_c$ \\
\hline
 50   & $130{\cdot}2$  & $8{\cdot}35\times10^{9}$ \\
 100  & $130{\cdot}2$  &           \\
 500  & $130{\cdot}9$  & $8{\cdot}15\times10^{9}$ \\
 1{,}000 & $130{\cdot}2$  &       \\
 2{,}000 & $130{\cdot}9$  & $8{\cdot}54\times10^{9}$ \\
 5{,}000 & $130{\cdot}9$ & $7{\cdot}79\times10^{9}$ \\
 10{,}000 & $3{\cdot}14$  & $3{\cdot}56\times10^{8}$\\
 50{,}000 & $3{\cdot}14$  &       \\
 100{,}000 & $3{\cdot}71$ & $5{\cdot}07\times10^{8}$ \\
 300{,}000 & $7{\cdot}47$ & $3{\cdot}43\times10^{8}$ \\
 1{,}000{,}000 & $17{\cdot}2$  &  \\
 3{,}000{,}000 & $18{\cdot}94$ & $3{\cdot}43\times10^{8}$ \\
 10{,}000{,}000 &             &   \\
 30{,}000{,}000 & $44{\cdot}6$ & $8{\cdot}94\times10^{8}$ \\
 300{,}000{,}000 &              & $8{\cdot}90$  \\
\hline
\end{tabular}
\end{table}

% ----------------------------------------------------------------
% Figures 47 and 48
% ----------------------------------------------------------------

\figurePlaceholder{48}{The dependence on $Q_1$ (for various assigned values
of~$T_1$) of the wave number~$a$ (in the unit $1/d$) of the disturbance at
which instability first sets in as convection (solid line) and as
overstability (for $\mathrm{p} = 0{\cdot}025$) (broken line).  It will be
observed that a discontinuous change in~$a$ occurs when (for increasing~$Q_1$)
the manner of instability changes from overstability to cellular convection.}

of the problem presents some very unexpected features.  These result from
the fact that the curve $R(a)$ defined by equations~\eqref{eq:5-59}
and~\eqref{eq:5-60} has, for certain ranges of the parameters~$Q_1$
and~$T_1$, \emph{two minima} (see Fig.~49); and the minimum giving the lower
Rayleigh number, when $Q_1$ is less than a certain value, gives the higher
Rayleigh number when $Q_1$ is larger than this value.  Thus, for
$T_1 = 10^{4}$ and $Q_1 = 80$, the two minima occur

\figurePlaceholder{49}{The dependence of the Rayleigh number~$R_c$ for the
onset of instability on the wave number when $T_1 = 10^{4}$ and $Q_1$ has the
values 40 (curve labelled~a), 100 (curve labelled~b), and 200 (curve
labelled~c).  The abscissa is related to the wave number~$a$ by
$x = a^{2}/\pi^{2}$.  Notice the occurrence of two minima and their relative
disposition.}

for $a = 18{\cdot}3$ and $3{\cdot}38$ where
$R = 4{\cdot}00\times10^{4}$ and $4{\cdot}78\times10^{4}$, respectively;
while, for $Q_1 = 100$ the two minima occur for $a = 18{\cdot}2$ and
$3{\cdot}37$, where $R = 3{\cdot}98\times10^{4}$ and
$3{\cdot}93\times10^{4}$, respectively.  Consequently, for $Q_1$ slightly
less than~100, the wave number of the cells which appear at marginal
stability will suddenly decrease from $a = 18{\cdot}2$ to $a = 3{\cdot}4$.
In other words, if we start with an initial situation in which
$T_1 = 10^{4}$ and no magnetic field is present, and gradually increase the
strength of the magnetic field, then at first the cells which appear at
marginal stability will be elongated; and when the magnetic field has
increased to a value corresponding to $Q_1 = 100$, cells of two very
different sizes will appear simultaneously: one set which will be highly
elongated and another set which will, relatively, be highly flattened.  As the
magnetic field increases further, only the flattened cells will persist.
%
% NOTE: The preceding sentence completes a gap in the sec51_52 fragment,
% which ended mid-sentence at "As the magnetic field increases".
% The phrasing is inferred from context; consult the original book (p.~207)
% for the exact wording.

%% ==================================================================
%% Section 53
%% ==================================================================
\section{The case when instability sets in as overstability}\label{sec:53}

We have seen that when liquid metals, such as mercury, are in rotation,
thermal instability sets in mostly as overstability. But in the presence of
a magnetic field, the principle of the exchange of stabilities is valid and
instability sets in as stationary convection. When rotation and magnetic
field are simultaneously present, the manner of the onset of instability
must, in general, depend in an extremely complicated way on the
relevant parameters $T$, $Q$, $p_1$ ($= \nu/\kappa$), and $p_2$ ($= \nu/\eta$). We shall seek
to unravel the nature of this dependence in the particular case of mercury.

Returning to equation~(43), and considering again the case of two free
boundaries adjoining a non-conducting medium, we can obtain the
required characteristic equation by inserting the solution $W = W_0 \sin \pi z$;
thus
\begin{align}
Ra^2(n^2{+}a^2{+}p_2\,\sigma)\bigl[(n^2{+}a^2{+}\sigma)(n^2{+}a^2{+}p_2\,\sigma){+}Qn^2\bigr]
\notag\\
= (n^2{+}a^2{+}p_1\,\sigma)\bigl\{(n^2{+}a^2)\bigl[(n^2{+}a^2{+}\sigma)(n^2{+}a^2{+}p_2\,\sigma){+}Qn^2\bigr]^2 +{} \notag\\
\qquad {}+ Tn^2(n^2{+}a^2{+}p_2\,\sigma)^2\bigr\},
\label{eq:5-62}
\end{align}
where it must be remembered that $\sigma$ can be complex. Letting
\begin{equation}
x = \frac{a^2}{\pi^2},\quad
i\sigma_i = \frac{\sigma}{\pi^2},\quad
R_* = \frac{R}{\pi^4},\quad
T_* = \frac{T}{\pi^4},\quad\text{and}\quad
Q_* = \frac{Q}{\pi^2},
\label{eq:5-63}
\end{equation}
we can rewrite equation~\eqref{eq:5-62} in the form
\begin{align}
R_* &= \frac{1+x}{x}\bigl(1{+}x{+}ip_1\,\sigma_i\bigr)\bigl(1{+}x{+}i\sigma_i\bigr)
+ Q_*\,\frac{1{+}x{+}ip_1\,\sigma_i}{1{+}x{+}ip_2\,\sigma_i}
\notag\\
&\quad{}+ \frac{T_*}{1{+}x}\;\frac{(1{+}x{+}ip_1\,\sigma_i)(1{+}x{+}ip_2\,\sigma_i)}
{\bigl(1{+}x{+}i\sigma_i\bigr)\bigl(1{+}x{+}ip_2\,\sigma_i\bigr)+Q_*}\,.
\label{eq:5-64}
\end{align}

Since our present interest is to specify the critical Rayleigh number
for the onset of instability via a state of purely oscillatory motions, it
will suffice to seek the conditions that equation~\eqref{eq:5-64} will allow solutions
for which $\sigma_i$ is real. Assuming then that this is the case, and equating
the real and the imaginary parts of equation~\eqref{eq:5-64}, we obtain the following
pair of equations for determining $R_*$ and $\sigma_i$ for given $x$, $Q_*$, and~$T_*$:
\begin{align}
R_* &= \frac{1+x}{x}\bigl[(1{+}x)^2 - p_1\,\sigma_i^2\bigr]
+ Q_*\,\frac{(1{+}x)^2{+}p_1\,p_2\,\sigma_i^2}{(1{+}x)^2{+}p_2^2\,\sigma_i^2}
\notag\\[6pt]
&\quad{}+ \frac{T_*}{1{+}x}\;\frac{\bigl[(1{+}x)^2{-}p_1\,\sigma_i^2{+}Q_*\bigr]\bigl[(1{+}x)^2{-}p_2\,p_1\,\sigma_i^2\bigr]
{+}(1{+}x)^2(p_1{+}p_2)(1{+}p_2)\sigma_i^2}
{\bigl[(1{+}x)^2{-}p_2\,\sigma_i^2{+}Q_*\bigr]^2{+}(1{+}x)^2(1{+}p_2)^2\sigma_i^4}
\label{eq:5-65}
\end{align}
and
\begin{align}
&\frac{T_*}{1{+}x}\;\frac{(p_1{-}1)(1{+}x)^2{+}(p_1{+}p_2)Q_*{+}p_2^2(p_1{-}1)\sigma_i^2}
{\bigl[(1{+}x)^2{-}p_2\,\sigma_i^2{+}Q_*\bigr]^2{+}(1{+}x)^2(1{+}p_2)^2\sigma_i^4}
\notag\\[6pt]
&\qquad{}+ Q_*\,\frac{p_1{-}p_2}{(1{+}x)^2{+}p_2^2\,\sigma_i^2}
+ 1 + p_1 = 0.
\label{eq:5-66}
\end{align}

For assigned values of $Q_*$ and $T_*$, equations~\eqref{eq:5-65} and~\eqref{eq:5-66} define $R_*$
as a function of $x$; the minimum of this function determines the critical
Rayleigh number for the onset of overstability (for the assigned $Q_*$ and
$T_*$). This minimum value of the Rayleigh number for the onset of over\-stability should be compared with the corresponding value for the onset
of convection as given by the formulae of \S\,52; the manner in which
instability will first set in will depend on which of the two Rayleigh
numbers is the smaller.

Of course, for certain ranges of the parameters, equation~\eqref{eq:5-66} will
not allow positive solutions for $\sigma_i^2$; for those ranges of the parameters,
overstability is not possible.

\subsection*{(a) An approximate solution applicable to liquid metals}

The equations giving $R_*$ and $\sigma_i$ simplify considerably when applied
to liquid metals such as mercury. These substances are characterized by
values of $p_2$ which are exceedingly small; thus for mercury at room
temperatures,
\begin{equation}
p_1 = \frac{1{\cdot}11}{4{\cdot}45}\times 10^{-1} = 0{\cdot}025;\qquad
p_2 = \frac{1{\cdot}11 \times 10^{-3}}{7{\cdot}6 \times 10^{4}} = 1{\cdot}5\times 10^{-7}.
\label{eq:5-67}
\end{equation}
Accordingly, in these cases, $p_2$ can be neglected in comparison with $p_1$
or unity. Moreover, it appears on examination that the terms in $p_2$ can
be neglected even when they occur multiplied by a large factor such as
$\sigma_i^2$. Therefore, suppressing all terms in equations~\eqref{eq:5-65}
and~\eqref{eq:5-66} which have $p_2$ as a factor, we obtain
\begin{equation}
R_* = \frac{1+x}{x}\bigl[(1{+}x)^2 - p_1\,\sigma_i^2\bigr]
+ Q_* + \frac{T_*}{1{+}x}\;\frac{(1{+}x)^2\bigl[(1{+}x)^2{+}Q_*\bigr]{+}\sigma_i^2}
{\bigl[(1{+}x)^2{+}Q_*\bigr]^2{+}(1{+}x)^2\sigma_i^2}\,,
\label{eq:5-68}
\end{equation}
and
\begin{equation}
(1{+}x)(1{+}p_1){+}Q_*\,\frac{p_1}{1{+}x} = \frac{T_*}{1{+}x}\;
\frac{(1{-}p_1)(1{+}x)^2 - p_1\,Q_*}
{\bigl[(1{+}x)^2{+}Q_*\bigr]^2{+}(1{+}x)^2\sigma_i^2}\,.
\label{eq:5-69}
\end{equation}
Equation~\eqref{eq:5-69} provides the following explicit formula for $\sigma_i^2$:
\begin{equation}
\sigma_i^2 = \frac{T_*}{1{+}x}\;
\frac{(1{+}x)^2(1{-}p_1) - p_1\,Q_*}{(1{+}x)(1{+}p_1)+p_1\,Q_*}
- \biggl[\frac{(1{+}x)^2{+}Q_*}{1{+}x}\biggr]^2.
\label{eq:5-70}
\end{equation}
Also, equations~\eqref{eq:5-68} and~\eqref{eq:5-69} can be combined to give
\begin{equation}
R_* = 2\,\frac{1+x}{x}\;\frac{(1{+}x)^2{+}Q_*}
{(1{+}x)^2(1{-}p_1)}\;\bigl[(1{+}x)^2{+}p_1\sigma_i^2\bigr].
\label{eq:5-71}
\end{equation}

\subsection*{(b) Numerical results for mercury}

Equations~\eqref{eq:5-70} and~\eqref{eq:5-71} have been used to determine the critical
Rayleigh numbers for the onset of overstability for various values of
$Q_*$ and $T_*$ for $p_1 = 0{\cdot}025$. The results of the calculations are summarized
in Table~\ref{tab:XIX}. The last column in this table gives the gyration frequency
of the oscillation in the marginal state in units of $\Omega$; according to equations~\eqref{eq:5-63} and the definition of~$T_*$:
\begin{equation}
p/\Omega = 2\sigma_i/\!\sqrt{T_*}\,.
\label{eq:5-72}
\end{equation}

\begin{table}[htbp]
\centering
\caption{Critical Rayleigh numbers and related constants for the onset of overstability
for the case $p_1 = 0{\cdot}025$}
\label{tab:XIX}
\medskip
{\small
\begin{tabular}{@{}r r r r r@{}}
\hline
$Q_*$ & $a_c$ & $R_{*c}$ & $R_c$ & $p/\Omega$ \\
\hline
\multicolumn{5}{c}{$T_* = 10^4$} \\[3pt]
10 & $4{\cdot}36$ & $33{\cdot}58$ & $7{\cdot}059\times10^3$ & $1{\cdot}0704$ \\
30 & $5{\cdot}35$ & $45{\cdot}35$ & $1{\cdot}388\times10^4$ & $0{\cdot}9110$ \\
100 & $6{\cdot}63$ & $90{\cdot}12$ & $3{\cdot}546\times10^4$ & $0{\cdot}6014$ \\
150 & $7{\cdot}30$ & $18{\cdot}44$ & $5{\cdot}019\times10^4$ & $0{\cdot}3688$ \\[6pt]
\multicolumn{5}{c}{$T_* = 10^5$} \\[3pt]
10 & $6{\cdot}03$ & $140{\cdot}8$ & $1{\cdot}284\times10^5$ & $0{\cdot}8905$ \\
20 & $6{\cdot}41$ & $139{\cdot}6$ & $1{\cdot}067\times10^5$ & $0{\cdot}8965$ \\
30 & $6{\cdot}66$ & $137{\cdot}8$ & $4{\cdot}039\times10^5$ & $0{\cdot}8803$ \\
40 & $6{\cdot}88$ & $133{\cdot}7$ & $2{\cdot}355\times10^4$ & $0{\cdot}7813$ \\
50 & $7{\cdot}03$ & $130{\cdot}4$ & $1{\cdot}499\times10^4$ & $0{\cdot}7615$ \\
70 & $7{\cdot}30$ & $117{\cdot}6$ & $1{\cdot}990\times10^4$ & $0{\cdot}7438$ \\
100 & $7{\cdot}66$ & $100{\cdot}9$ & $4{\cdot}317\times10^4$ & $0{\cdot}6894$ \\
200 & $8{\cdot}45$ & $94{\cdot}38$ & $7{\cdot}139\times10^4$ & $0{\cdot}5099$ \\
400 & $9{\cdot}48$ & $74{\cdot}53$ & $1{\cdot}278\times10^5$ & $0{\cdot}4587$ \\
300 & $9{\cdot}88$ & $61{\cdot}95$ & $1{\cdot}557\times10^5$ & $0{\cdot}3918$ \\[6pt]
\multicolumn{5}{c}{$T_* = 10^6$} \\[3pt]
10 & $8{\cdot}35$ & $341{\cdot}1$ & $4{\cdot}690\times10^5$ & $0{\cdot}6842$ \\
30 & $8{\cdot}89$ & $311{\cdot}7$ & $4{\cdot}533\times10^5$ & $0{\cdot}6434$ \\
100 & $9{\cdot}78$ & $293{\cdot}6$ & $7{\cdot}073\times10^5$ & $0{\cdot}6520$ \\
300 & $11{\cdot}34$ & $196{\cdot}5$ & $1{\cdot}893\times10^5$ & $0{\cdot}4760$ \\
1{,}000 & $12{\cdot}64$ & $202{\cdot}3$ & $3{\cdot}371\times10^5$ & $0{\cdot}4050$ \\
1{,}600 & $13{\cdot}62$ & $179{\cdot}3$ & $4{\cdot}898\times10^5$ & $0{\cdot}3406$ \\[6pt]
\multicolumn{5}{c}{$T_* = 10^7$} \\[3pt]
10 & $17{\cdot}4$ & $1739$ & $6{\cdot}000\times10^6$ & $0{\cdot}3458$ \\
100 & $18{\cdot}6$ & $1674$ & $6{\cdot}420\times10^6$ & $0{\cdot}3348$ \\
1{,}000 & $20{\cdot}9$ & $1407$ & $1{\cdot}002\times10^6$ & $0{\cdot}3136$ \\
10{,}000 & $25{\cdot}9$ & $1133$ & $3{\cdot}642\times10^6$ & $0{\cdot}2246$ \\
15{,}000 & $28{\cdot}4$ & $1024$ & $5{\cdot}348\times10^6$ & $0{\cdot}2048$ \\[6pt]
\multicolumn{5}{c}{$T_* = 10^8$} \\[3pt]
100 & $37{\cdot}7$ & $8097$ & $1{\cdot}354\times10^7$ & $0{\cdot}1619$ \\
1{,}000 & $38{\cdot}4$ & $7933$ & $1{\cdot}305\times10^7$ & $0{\cdot}1591$ \\
10{,}000 & $41{\cdot}9$ & $7289$ & $1{\cdot}720\times10^7$ & $0{\cdot}1446$ \\
40{,}000 & $46{\cdot}5$ & $6453$ & $2{\cdot}786\times10^7$ & $0{\cdot}1491$ \\
100{,}000 & $50{\cdot}1$ & $5841$ & $4{\cdot}376\times10^7$ & $0{\cdot}1168$ \\
130{,}000 & $51{\cdot}3$ & $5643$ & $5{\cdot}421\times10^7$ & $0{\cdot}1119$ \\
\hline
\end{tabular}
}
\end{table}

The critical Rayleigh numbers for the onset of overstability and
stationary convection are exhibited in Fig.~47 as functions of $Q_*$ for
various assigned values of $T_*$; the wave numbers of the associated disturbances are similarly exhibited in Fig.~48. We can directly read from
Fig.~47 the manner of the onset of instability for the various values of
$T_*$ for which the calculations have been made.

The phenomenon described in \S\,52 which gives rise to the peculiar,
non-monotone behaviour of the $(R, Q)$-curves when instability sets in
as stationary convection does not occur when allowance is made for the
onset of overstability for a value of $p_1$ as small as $0{\cdot}025$. However, the
curves do show a discontinuous behaviour though of a much less pro\-nounced character. From Fig.~47 it appears that for the cases illustrated,
the transition from overstability to convection occurs at about the place
where the convection curve passes through its minimum; and from Fig.~48 it is seen that this transition is accompanied by a substantial dis\-continuity in the wave number of the disturbance which manifests itself
at marginal stability. This discontinuity is in the sense that the convection cells (for increasing $Q_*$) suddenly get very much widened.

%% ==================================================================
%% Section 54
%% ==================================================================
\section{Experiments on the onset of thermal instability in the
presence of rotation and magnetic field}\label{sec:54}

Experiments on the onset of thermal instability in the presence of
rotation and magnetic field have been successfully carried out by
Nakagawa.

In Nakagawa's experiments, a reconditioned electromagnet of a $32\tfrac{1}{2}$-in.\
cyclotron of the University of Chicago was used to provide a uniform
vertical magnetic field of variable strength. A Pyrex glass cylinder
containing mercury mounted on levelled non-magnetic ball-bearings
between the pole pieces of the electromagnet was rotated in the magnetic
field by a constant speed electric motor (see Fig.~50 which shows the
arrangement used in the experiments described in \S\,(b) below). The
angular velocity of the rotation was measured by an optical method using
electrical timing devices; the accuracy attained in the time measurements was $10^{-3}$~sec.

A direct-current electric heater placed at the bottom of the container,
below a stainless-steel plate, provided the means of heating. The
temperature gradient which was established was measured by a nine-couple copper--constantan thermopile immersed at two fixed levels in
the mercury. The average temperature of the mercury in any particular
experiment was also measured by a thermocouple, one of whose junctions
was placed at the middle of the mercury layer; the values of the various
coefficients (like $\nu$, $\kappa$, etc.)\ at this average temperature were used in
evaluating the parameters such as $R$, $T$, and~$Q$.

The e.m.f.\ of the thermopile and the thermocouple were amplified
and recorded automatically and continuously. The accuracy in the
thermopile measurements was $\pm 0{\cdot}001^\circ$~C, while the average temperatures measured by the thermocouple were accurate to $\pm 0{\cdot}01^\circ$~C.

The electrical connexions to the various components in the rotating
part of the assembly were made through a system of mercury filled
troughs with copper and constantan contacts.

A constant circulation of cooled nitrogen helped to maintain the
conditions at the top surface of the mercury constant and uniform.

\figurePlaceholder{50}{A schematic diagram of the experimental arrangement:
$A$, bakelite cylinder; $B$, stainless-steel plate; $C$, electric heater;
$D$, non-magnetic ball-bearing; $E$, stainless-steel rod; $F$, mercury trough;
$M$, front-surface mirror; $S$, rotary shutter; $T$, camera
(\textit{Proc.\ Roy.\ Soc.\ (London)} A, \textbf{249}, 140 (1958)).}

\subsection*{(a) The results on the critical Rayleigh number and on the manner of
the onset of instability}

All the experiments were carried out with a layer of mercury 3~cm in
depth rotated with a constant angular velocity of 5~rev/min. The corresponding value of $T_*$ varied between $7{\cdot}5\times 10^4$ and $8{\cdot}5\times 10^6$ depending on
the average temperature of the mercury. Experiments were performed
for values of $Q_*$ in the range 10 to $2\times 10^4$ by using magnetic fields of
strength varying from 125 to 5{,}000~gauss.

The critical Rayleigh number for the onset of instability, for a constant
speed of rotation and strength of magnetic field, was determined by
measuring the steady temperature gradients which get established for
various rates of heating, and using the Schmidt--Milverton principle;
and the manner of the onset of instability was ascertained in each case
by inspecting the temperature records. The onset of overstability could
always be distinguished by the nearly pure sinusoidal oscillations
exhibited by the temperature records. A typical example of such a
record is shown in Fig.~51.

\figurePlaceholder{51}{A time record of the adverse temperature gradient for
mercury; $d = 3$~cm; $\Omega = 5$~rev/min; $H = 125$~G; $Q_* = 1{\cdot}01\times 10^1$;
$T = 7{\cdot}90\times 10^4$.}

In Fig.~52 the results of the experiments together with the theoretically
expected relation for $T_* = 10^6$, according to the calculations of \S\S\,52
and~53, are shown. Since the theoretical results have been derived for two
free boundaries and refer, moreover, to a value of $T_*$ somewhat different
from the average value to which the experiments correspond, a quantitative agreement between the experiments and the theory is not to be
expected. All that can be expected is an agreement in the general
features; and this is certainly present. In particular, the characteristic
discontinuity in the dependence of the Rayleigh number on $Q_*$ with the
accompanying change in the manner of the onset of instability is clearly
demonstrated.

The periods $2\pi/p$ of the overstable oscillations were estimated by
counting the number of oscillations in the temperature records which
occur in known intervals of time. In Fig.~53 these results are shown
together with the theoretically expected relations for $T_* = 10^5$ and~$10^6$.
Again the general agreement is satisfactory, though as we have explained,
quantitative agreement is not to be expected.

\figurePlaceholder{52}{A summary of the experimental and the theoretical results on the
critical Rayleigh numbers for instability. The curves labelled $T_* = 10^6$ (convection), $T_* = 10^6$ (overstability), and $T_* = 0$ (convection) represent the
theoretically derived relations. The value of $T_*$ appropriate for the experimental points represented by the solid and shaded triangles is $7{\cdot}78\times 10^5$.}

\figurePlaceholder{53}{A comparison of the observed periods of overstable oscillations with
the theoretical values (the full line curves). The value of $T_*$ appropriate for
the experimental results is $(6{\cdot}05\pm 0{\cdot}07)\times 10^5$.}

\subsection*{(b) Optical observations and the discontinuous variation of the cell
dimensions with the strength of the magnetic field}

For optical observations, some slight modifications in the arrangement
described earlier were necessary. It is the arrangement used in these
experiments that is illustrated in Fig.~50.

Mercury was contained in a Bakelite cylinder 24~cm in diameter and
4~cm in height. The bottom of the cylinder $A$ was sealed by a stainless-steel plate $B$ with an O-ring. The steel plate was carefully machined to
eliminate surface irregularities which might disturb the ensuing convection patterns; freedom from such disturbances was found to be essential
for successful optical observations.

Uniform heating from below was provided by a non-inductively
wound electric heater $C$ placed below the stainless-steel plate, and
the electric heater was mounted, between the pole pieces of the electromagnet, on non-magnetic ball-bearing~$D$. And this assembly was rotated
in the magnetic field at 5~rev/min as in the previous experiments.

The optical part of the apparatus consisted of a mirror $M$ mounted $45^\circ$
to the vertical, the rotoscope, a rotary shutter~$S$, and a single lens reflex
camera~$T$. The mirror was silvered on the front surface; and dimming
of the mirror (by condensation of water-vapour during the experiments)
was prevented by heat from an infra-red lamp.

A stationary image was secured at the camera by rotating the rotoscope with an angular velocity exactly one-half that of the apparatus.
By a suitable sequence of openings and closings of the shutter, it was
possible to obtain photographs every 15~sec.

As tracers of the motion, sand particles of diameter approximately
$\frac{1}{4}$~mm were used. These particles floating freely on the surface, and
following the motion of the mercury, enabled the emerging pattern of
convection to be photographed. A thin layer of distilled water on the
mercury surface eliminated much of the difficulty arising from surface
oxidation.

The experiments were again carried out on a layer of mercury 3~cm in
depth and an angular velocity of rotation of 5~rev/min. And as ascertained by the experiments in~\S\,(a), the rate of heating for each value of $Q_*$
was adjusted so that the conditions were just marginal. By repeating
the experiments at least twice for each value of $Q_*$, Nakagawa took a
total of 300 streak photographs. Measurements were made on prints of
these photographs on $8\times 10$-in.\ paper. Distances between the centres
of adjacent cells were measured and in the manner described in Chapter~IV, \S\,48, the theoretical values to which these distances should correspond were inferred. In Fig.~54 the results of the measurements are
compared with the expected theoretical relation for $T_* = 10^6$. While,
for reasons we have already explained, a quantitative agreement is not
to be expected, the agreement in the general features of the relationship

\figurePlaceholder{54}{A summary of the experimental and theoretical results on the sizes
of the cells manifested at marginal stability. The theoretical relations for
$T_* = 10^6$ and $T_* = 0$ are represented by the broken lines. Solid circles
represent the experimental results for $T_* = 7{\cdot}30\times 10^5$; and the heavy solid lines
are fitted to the experimental data. The lower branches of the broken line
and the heavy solid line represent the case when the instability sets in as
overstability. Both lines on the right-hand side of the figure represent the
case when the instability sets in as steady convection.}

\noindent
is satisfactory. In particular, the predicted enlargement of the cell
dimension at a critical field strength is strikingly confirmed.

In Fig.~55 some typical examples from Nakagawa's collection of
streak photographs are illustrated. They illustrate the characteristic
dependence of the size of the cells on the value of $Q_*$. The discontinuous
change in the size of the cells which occurs for increasing field strength is
clearly demonstrated by this sequence of photographs.

\figurePlaceholder{55}{Examples of streak photographs of the convective motion on the surface of a layer
of mercury of depth 3~cm, rotating at 5~rev/min, and subject to impressed external magnetic
fields of various strengths: (a) $H = 125$~G, $Q_* = 9{\cdot}60$, and $T_* = 6{\cdot}95\cdot 10^5$;
(b) $H = 750$~G, $Q_* = 3{\cdot}52\cdot 10^3$, and $T_* = 7{\cdot}30\cdot 10^5$;
(c) $H = 1{,}000$~G, $Q_* = 6{\cdot}25\cdot 10^3$, and $T_* = 7{\cdot}31\cdot 10^5$;
(d) $H = 3{,}000$~G, $Q_* = 5{\cdot}61\cdot 10^4$, and $T_* = 7{\cdot}49\cdot 10^5$.
The sudden enlargement of the cell size as we pass from $H = 750$~G to $H = 1{,}000$~G is apparent.}

%% ==================================================================
%% BIBLIOGRAPHICAL NOTES
%% ==================================================================
\section*{BIBLIOGRAPHICAL NOTES}
\addcontentsline{toc}{section}{Bibliographical Notes}

\paragraph{\S\,50.} Hydromagnetic waves in a rotating fluid have been considered by:
\begin{enumerate}
\item B.~\textsc{Lehnert}, `Magnetohydrodynamic waves under the action of the
Coriolis force, Part~I', \textit{Astrophys.\ J.}\ \textbf{119}, 647--54 (1954);
\textbf{Part~II}, \textit{ibid.}\ \textbf{121}, 481--90 (1955).
\end{enumerate}
See also:
\begin{enumerate}\setcounter{enumi}{1}
\item S.~\textsc{Chandrasekhar}, `The gravitational instability of an infinite homogeneous medium when Coriolis force is acting and a magnetic field is present',
\textit{ibid.}\ \textbf{119}, 7--9 (1954).
\item B.~\textsc{Lehnert}, `The decay of magneto-turbulence in the presence of a
magnetic field and Coriolis force', \textit{Quart.\ Appl.\ Math.}\ \textbf{12}, 321--41 (1955).
\end{enumerate}

\paragraph{\S\S\,51--53.} The analyses in these sections are based on:
\begin{enumerate}\setcounter{enumi}{3}
\item S.~\textsc{Chandrasekhar}, `The instability of a layer of fluid heated below and
subject to the simultaneous action of a magnetic field and rotation.~I',
\textit{Proc.\ Roy.\ Soc.\ (London)}~A, \textbf{225}, 173--84 (1954); II, \textit{ibid.}\ \textbf{237}, 476--84 (1956).
\end{enumerate}

The fact that the curve $R(x)$ defined by equations~(59) and~(60) has two minima
for certain ranges of the parameters $Q$ and $T$, was first observed by Donna Elbert.

\paragraph{\S\,54.} The references for the experimental work described in this section are:
\begin{enumerate}\setcounter{enumi}{4}
\item Y.~\textsc{Nakagawa}, `Experiments on the instability of a layer of mercury
heated from below and subject to the simultaneous action of a magnetic
field and rotation.~I', \textit{Proc.\ Roy.\ Soc.\ (London)}~A, \textbf{242}, 81--88 (1957);
II, \textit{ibid.}\ \textbf{249}, 138--45 (1959).
\end{enumerate}

Extending the experiments described in reference~5, Nakagawa has since demonstrated the occurrence of instability as stationary convection after the instability
has already manifested itself as overstability; and he has, moreover, observed
the predicted decrease of the Rayleigh number, with increasing strength of the
magnetic field, over a part of the convective branch (see Fig.~52). [For the corresponding results for the case when rotation alone is present, see pp.~142 and~143.]
